\section{Experimental protocol}
\label{sec:protocol}

\textbf{Graphs.} Synthetic ($n=4000$ unless stated): Erd\H{o}s--R\'enyi (ER) and random regular with mean degree $d\in\{4,\dots,256\}$; stochastic block model (SBM, 40 equal blocks, $p_{\rm in}/p_{\rm out}=30$); Watts--Strogatz (WS, rewiring $0.05$; strong local clustering); Barab\'asi--Albert (BA, $m\in\{3,\dots,24\}$; heavy-tailed degrees). Real: Stanford3 (Facebook100~\citep{traud2012social}; Network Repository~\citep{rossi2015network}; $n=11{,}586$, $568{,}309$ edges, mean/median/max degree $98/73/1172$, mean local clustering $0.24$) and ego-Facebook (SNAP~\citep{mcauley2012learning}; $n=4{,}039$, $88{,}234$ edges, mean/median/max degree $44/25/1045$, clustering $0.61$). Both are connected. Data versions and checksums are in Appendix~\ref{app:repro}. Clusterability differs sharply: the Louvain partition into $\approx30$ clusters cuts $18\%$ of ego-Facebook's edges but $58\%$ of Stanford3's.

\textbf{Potential outcomes.} Unit-level draws are made once per world: baseline $a_i\sim N(0,1)$, direct effect $\beta_i=1+0.5\xi_i$, spillover coefficient $\gamma_i=\gamma(1+0.5\zeta_i)$, $\xi,\zeta\sim N(0,1)$ (so $\bar\beta\approx1$, $\bar\gamma\approx\gamma$, default $\gamma=1$). There is no extra observation noise; baselines and effect heterogeneity act as design-based noise. Mechanisms:
\textbf{(i)~one-hop fraction}, $Y_i=a_i+\beta_iz_i+\gamma_ig(s_i)$ with $g(s)=s^k$ (linear $k=1$; concave $k\in\{\frac14,\frac12\}$; convex $k\in\{2,4\}$) or a threshold response $g(s)\propto\sigma((s-\frac12)/0.08)$ rescaled to $g(0)=0,g(1)=1$ so that $\tau\approx\bar\beta+\bar\gamma$ for every shape; \textbf{(ii)~multi-hop} (linear-in-means equilibrium) $Y=(I-\lambda W)^{-1}(a+\beta z)$, $\lambda\in[0,0.8]$, $\tau=\bar\beta/(1-\lambda)$ in expectation; \textbf{(iii)~count}: $Y_i=a_i+\beta_iz_i+\gamma_ic_i/\bar d$; \textbf{(iv)~degree-heterogeneous} spillover: (i) with $\gamma_i\propto e^{\mp0.7\ell_i}$ ($\ell_i$ standardized log degree; ``inv'': low-degree units more susceptible; ``pos'': high-degree units), normalized to mean $\gamma$. The true $\tau$ is computed exactly from $Y(\mathbf1)-Y(\mathbf0)$.

\textbf{Experiments and worlds.} Seven experiments (E1--E7; scripts in Appendix~\ref{app:repro}) vary degree, $n$ and noise (E1); response shape (E2); neighbourhood, heterogeneity and exposure definition (E3); spillover strength (E4); treated fraction and number of clusters (E5); inference (E6); a mis-measured graph (E7). Each setting uses $6$ independent \emph{worlds} (a fresh synthetic graph and unit parameters, or fresh unit parameters on the fixed real graph) with $R=300$ assignments each ($400$ in E6, $40$ for \textsc{gbm}). Hypotheses H1--H6 (error $\propto$ leverage; persistent extrapolation bias; multi-hop omission bias; heterogeneity bias; clustering helps iff the graph clusters; regression beats DiM only above a spillover threshold) were fixed before the final runs. A three-world pilot of E1/E3 set the run scale and exposed a crash (assignments with an empty arm under few clusters), fixed in the released code; \textsc{wls-s} and \textsc{strat-s} were added after the pilot heterogeneity results and are labelled as such.
